\chapter{Reading and using the library}\label{chap:foundations}
\section{Scope and a first working program}
\api{nurbspath} is a dependency-free, header-only C++20 library. Its geometry
types use a scalar template parameter constrained by \api{std::floating_point};
\api{float}, \api{double}, and \api{long double} are valid choices. Include
\api{nurbspath/nurbspath.hpp} for the complete API or a focused public header
for a smaller translation unit.

This manual describes the source snapshot printed on the title page. Chapters
\ref{chap:entities}--\ref{chap:output} explain the public operations by purpose;
Appendix~\ref{chap:reference} records individual overloads, defaults, result
types, and public data members extracted from those headers. Internal
\api{nurbspath::detail} helpers are implementation, not a public contract.

\begin{lstlisting}
#include <nurbspath/config.hpp>
#include <nurbspath/nurbspath.hpp>
#include <iomanip>
#include <iostream>

int main() {
    namespace np = nurbspath;
    using REAL = double;
    const np::circle2<REAL> circle({0.0, 0.0}, 1.0);
    const np::ray2<REAL> ray({-3.0, 0.0}, {2.0, 0.0});
    const auto result = np::intersect(ray, circle);
    std::cout << std::scientific << std::setprecision(4);
    for (const auto& hit : result.points) {
        std::cout << "s=" << hit.s << " point=" << hit.point << '\n';
    }
}
\end{lstlisting}
\begin{lstlisting}[style=plain]
s=1.0000e+00 point=-1.0000e+00 0.0000e+00
s=2.0000e+00 point=1.0000e+00 0.0000e+00
\end{lstlisting}
The ray direction has length two, so the native parameter is half the travelled
world distance. Compile this program directly with the source include directory:
\begin{lstlisting}[style=plain]
c++ -std=c++20 -I../include first_query.cpp -o first_query
\end{lstlisting}
This command is shown relative to the manual directory. The manual's own
reproducible reference programs and build instructions are in
Appendix~\ref{chap:reproduction}.

\section{Two independent Cartesian worlds}
\begin{tabularx}{\linewidth}{@{}lYY@{}}
\toprule
Concept & Independent 2D world & Shared 3D world\\
\midrule
Position & \api{point2<REAL>} & \api{point3<REAL>}\\
Offset or direction & \api{vector2<REAL>} & \api{vector3<REAL>}\\
Forward line & \api{ray2<REAL>} & \api{ray3<REAL>}\\
Round surface & \api{circle2<REAL>} & \api{sphere3<REAL>}\\
Rational path & \api{nurbs_spline2<REAL>} & \api{nurbs_spline3<REAL>}\\
Array-backed path & \api{nurbs_arr_spline2}\newline\api{<REAL>} & \api{nurbs_arr_spline3}\newline\api{<REAL>}\\
Plane & No 2D plane entity & \api{plane3<REAL>}\\
\bottomrule
\end{tabularx}

A point represents a location; a vector represents an offset or direction.
Point subtraction gives a vector, point plus vector gives a point, and vector
arithmetic remains vector arithmetic. A 2D entity does not implicitly become
a 3D entity. The explicit bridge is \api{project(plane, entity2)}, described in
Chapter~\ref{chap:entities}. Queries on 2D entities compute entirely in 2D.

\section{Units, parameters, and tolerances}
The library does not prescribe metres, millimetres, or another length unit.
Choose one consistent world unit. Angles are in radians. Display widths and
text sizes are separate consumer-selected display units.

\begin{tabularx}{\linewidth}{@{}lYY@{}}
\toprule
Entity & Parameter domain & Meaning\\
\midrule
Ray & $s\geq0$ for forward queries & $O+s\vect d$; direction is unnormalised.\\
Spline & $[t_p,t_C]$ & Active knot interval; $C$ is the number of controls.\\
Circle & $u\in[0,2\pi)$ & Counterclockwise angle from positive $x$.\\
Sphere & $u\in[0,2\pi)$, $v\in[-\pi/2,\pi/2]$ & Longitude and latitude.\\
Plane & $(u,v)\in\Real^2$ & Orthonormal local coordinates.\\
\bottomrule
\end{tabularx}

The ray's evaluation formula can be evaluated for negative $s$, but intersection
queries treat it as a forward half-line. Spline parameters remain native:
interpolation preserves the supplied stations, and plane embedding preserves
the same $s$ values. A spline parameter is not automatically a physical arc
length everywhere between interpolation stations.

\begin{manualnote}[What a tolerance controls]
Definition tolerances validate entity construction and parameter boundaries.
\api{numerical_settings<REAL>} separately controls function residuals,
parameter intervals, iteration counts, and sampling density. These quantities
have different units. Sampling-based intersections, spline distances, and arc
length estimates depend on their chosen sample counts; a finite sampling run
does not prove a continuous global result for every possible curve.
\end{manualnote}

\section{Values, shared ownership, and checked editing}
Construct entities as ordinary values for simple ownership. The
\api{make_*} functions in \api{creators.hpp} mirror geometry constructors and
return \api{std::shared_ptr} objects. They preserve constructor validation and
exceptions. They do not introduce a common geometry base class.

Ordinary spline definitions are private and use checked, atomic setters.
Failed edits leave the old definition intact. Array-backed alternatives
deliberately expose a \api{std::valarray<REAL>}; direct replacement is allowed,
so geometry operations validate the current storage and do not rely on stale
endpoint caches. References returned by getters follow their documented
invalidation rules.

\section{How the worked figures were made}
Simple vector, plane, and circle constructions are drawn directly in TikZ.
Curves, analytic derivatives, interpolation points, intersections, and closest
points use data generated by \api{pictures/generate_reference.cpp}, which
includes the library itself. The generator checks its results before writing
the coordinate files. The stored plot coordinates and displayed output use
scientific notation with at most four decimal places. The calculations are
performed in full \api{double} precision before display rounding.

The figures distinguish a mathematical curve from its drawn approximation:
PGFPlots connects library-evaluated points. Its smooth appearance is not an
extra numerical guarantee. This manual requires no image downloads, external
plotting service, shell escape, or TikZ externalisation.
